\begin{claim}{D.2.30}{Proved}{The remaining once-suspended coefficient is realized through dimension ten}
Put $B=K(\FF_2,2)$, $A=Q\Sigma B$, $Y=\Sigma^2B$,
$Z=\fib(Y\to K(\FF_2,4))$, and $W=\Omega Z$.
Use the common class $\vartheta_7\in H^7W$ of D.2.28.
There is an actual ordinary map on a two-local ten-skeleton
of $A$,
\[
 g_{10}:K_{10}\longrightarrow W,
 \qquad
 g_{10}^*\vartheta_7=
       (\alpha_3\operatorname{Sq}^1\alpha_3)|_{K_{10}}\ne0.
 \tag{D.2.30a}
\]
Here the finite cellular model is chosen through dimension eleven
before cancelling odd-unit cell pairs, as in D.2.19--D.2.20.
The map can be pulled back to a corresponding ordinary finite
skeleton by cellular approximation and localization.

This realizes the value $\epsilon=1$ left in D.2.28 by one
coordinated choice of extensions. It does not extend the map
over the eleventh skeleton or over all of $A$, and it gives
no assertion about subsequent nullification stages.
\end{claim}
\begin{proof}
All homology and cohomology below have coefficients $\FF_2$,
except for the displayed homotopy groups and Moore spaces.
Write $h$ for the ordinary mod-two Hurewicz homomorphism.
Write $F=\Omega W$. Its classes $d,e,f$ in degree six are
\[
 d=b^2,\qquad e=b^2+\{a,b\},\qquad f=p_3,
 \qquad r=e+f.
\]
We use the same letters for their homology suspensions in
degree seven of $W$. Let $R_8\in H_8W$ be the suspension
of $R=\{a,c\}\in H_7F$.

\medskip\noindent
\textbf{The actual target groups and the missing Hurewicz direction.}
The target calculations D.2.18--D.2.19 and D.2.22 give
\[
 \begin{aligned}
 \pi_7W&=\ZZ/4\{\rho\},& h(\rho)&=r,\\
 \pi_8W&=\ZZ/2\{\lambda\},& h(\lambda)&=R_8,\\
 \pi_9W&\cong(\ZZ/2)^2,&\ker h
       &=\ZZ/2\{w_9\}.
 \end{aligned}
 \tag{D.2.30b}
\]
The last assertion uses the actual image calculation in
\textup{(D.2.22t19)}. The two actual eight-spheres in $F$
have Hurewicz images $A_4^2$ and $J_2$.
Homology suspension kills $A_4^2$ and retains $J_2$;
the path-fibration calculation of D.2.27 identifies
$H_9W$ with $H_8F/\langle A_4^2\rangle$.
Thus $w_9$ is the adjoint of the actual sphere with
Hurewicz image $A_4^2$. In particular, Hurewicz injectivity
for $F$ has not been transported incorrectly to $W$.

Let $a_5\in H^5W$ be the bottom class. In the actual
mapping cone $C_{w_9}$ it has a unique extension
$\widehat a_5$, and
\[
 a_5^2=0\text{ in }H^{10}W,
 \qquad
 \widehat a_5^{\,2}=\text{the nonzero relative top class}
       \text{ in }H^{10}C_{w_9}.
 \tag{D.2.30c}
\]
Here is the precise detector. In a minimal Adams--Hilton
model for the four-connected space $W$, the bottom
generator has degree four and represents $A_4$.
The adjoint of $w_9$ represents $A_4^2$ in
$H_8\Omega W=H_8F$. Attaching the ten-cell adds a
degree-nine generator whose differential represents this
class. Subtracting an old chain makes that differential
exactly $A_4^2$. Its quadratic part is the reduced diagonal
of the new top class, by the actual cell-attachment and
cup-product comparison in
\cite[\S3.1, p.~300]{FGHHalperinLemaire1995}.
This gives its square coefficient one.
The old square $a_5^2$ is zero: a nonzero old coefficient
would already make $A_4^2$ a boundary. In this degree no
other word can cancel it, since positive generators begin
in degree four and the bottom space is one-dimensional.
The nonzero established class $A_4^2\in H_8F$ excludes
that possibility. Thus the new square is precisely the
relative top class, proving \textup{(D.2.30c)}.

\medskip\noindent
\textbf{A coordinated ordinary map on the eight-skeleton.}
Use the finite two-local cellular models of
D.2.19--D.2.20 and put $Q_i=K_i/K_4$.
Choose the model through dimension eleven before
cancelling odd-unit pairs. Its mod-two cellular
differential is then zero in the dimensions used below.
Cancellation is performed on the actual two-local
cellular approximation, with higher attachments
transported along the equivalence; no unstable map is
replaced merely by its stable homology matrix.
The target $W$ is two-local, so the original mapping
problem is recovered by localization and cellular
approximation.

The actual eight-stage decomposition
\textup{(D.2.20s1)} is
\[
 Q_8\simeq D'\vee M(\ZZ/2,7)_F\vee M(\ZZ/4,7)_J
                \vee C_1^8\vee S^8_s.
 \tag{D.2.30d}
\]
Here $D'$ is the unit block, $J$ is the seven-sphere
with mod-two label $Q^4A_0+A_0B_0$,
\[
 C_1^8=\cofib\bigl((2,\eta_6):S^6\vee S^7\to S^6\bigr),
 \qquad s=Q^5A_0-A_0C_5.
\]
The top eight-cell of $C_1^8$ is $A_0C_5$.
The top cell of $M(\ZZ/4,7)_J$ is the slide
$B_0^2-A_0C_5$. All these comparisons are actual:
the quotients are four-connected and have dimension
at most eight, within the ordinary realization range
of the stable cell slides.

Set the map to zero on $D'$ and $M(\ZZ/2,7)_F$.
Send the bottom sphere of $M(\ZZ/4,7)_J$ to $\rho$.
It extends across the top cell because $4\rho=0$.
Any extension sends its top mod-two class to a multiple
of $R_8$: that source class has zero Bockstein and
zero second square, and the complete target table
in degree eight has simultaneous kernel exactly
$\langle R_8\rangle$.
Changing the extension by the actual sphere $\lambda$
toggles this coefficient. Choose the top value zero.

On $C_1^8$, send the six- and seven-skeleton to zero
and send the top eight-cell, through its pinch, to
$\lambda$. Send $s$ to $-\lambda$.
These ordinary maps define an actual $g_8:Q_8\to W$.
In the original source labels their values are
\[
 \begin{aligned}
 (g_8)_*(Q^4A_0)&=0,&(g_8)_*(A_0B_0)&=r,\\
 (g_8)_*(A_0C_5)&=R_8,&(g_8)_*(Q^5A_0)&=0,\\
 (g_8)_*(B_0^2)=R_8.
 \end{aligned}
 \tag{D.2.30e}
\]
The last value uses the stated slide of the Moore top.
The unit degree-seven classes have value zero.
Since $\vartheta_7$ is dual to the $e$ coordinate,
the complete source degree-seven basis of D.2.28 gives
\[
 g_8^*\vartheta_7=
          (\alpha_3\operatorname{Sq}^1\alpha_3)|_{Q_8},
 \qquad g_8^*a_5=0.
 \tag{D.2.30f}
\]
For the first formula, the notation on the right means
the unique quotient class pulling back to that class
on $K_8$. Collapsing $K_4$ is an isomorphism on
cohomology in degree seven.

\medskip\noindent
\textbf{Every ninth and tenth attachment extends.}
Each nine-cell has zero mod-two cellular boundary.
Its attaching map has zero mod-two Hurewicz in $Q_8$,
so its composite with $g_8$ has zero Hurewicz in $W$.
The middle line of \textup{(D.2.30b)} is injective.
All these composites are therefore null, and choices
of their nullhomotopies give an actual $g_9:Q_9\to W$.

The same cellular argument makes every ten-cell
composite belong to $\{0,w_9\}$.
Suppose one were $w_9$. Its actual attachment supplies
a map of cones
\[
 Q_9\cup e^{10}\longrightarrow C_{w_9}
 \tag{D.2.30g}
\]
whose relative top-cell degree is one. The source
ten-cell is a nonzero mod-two relative cohomology
coordinate, since its mod-two cellular boundary is
zero. Hence the relative top class in
\textup{(D.2.30c)} pulls back nontrivially.
On the other hand, $g_9^*a_5=0$ by
\textup{(D.2.30f)}, and attaching a ten-cell does not
change degree-five cohomology. The pullback of
$\widehat a_5$ in \textup{(D.2.30g)} is therefore
zero, so its square is zero. This contradicts
\textup{(D.2.30c)}.

Every tenth-cell obstruction is consequently zero.
The cell extensions attach to the same fixed $Q_9$,
so they assemble to an ordinary map $Q_{10}\to W$.
Its composite with $K_{10}\to Q_{10}$ is the map
in \textup{(D.2.30a)} and retains all the values
in \textup{(D.2.30e)}.
The next attaching group is $\pi_{10}W=\pi_9F$,
where D.2.22 has an integral-Hurewicz kernel and
an additional direction detected by the order-eight
homology class. This proof does not remove those
eleventh-cell obstructions.
\end{proof}
